\chapter{چک‌لیست آمادگی و پاسخ به منتقدان}
\label{app:checklist}

گذار دموکراتیک یک پروژه‌ی مهندسی اجتماعی بزرگ است که به آمادگی قبلی نیاز دارد.

\section{چک‌لیست آمادگی برای جریان‌های سیاسی}

هر حزب یا ائتلافی که قصد شرکت در مجلس مهستان را دارد باید از هم‌اکنون این موارد را آماده کند:

\begin{itemize}
    \item[$\square$] \textbf{پیش‌نویس قانون اساسی موقت:} ارائه‌ی دیدگاه خود درباره‌ی ساختار گذار.
    \item[$\square$] \textbf{برنامه‌ی ۱۰۰ روز اول:} اقدامات فوری در اقتصاد، امنیت و روابط خارجی.
    \item[$\square$] \textbf{فهرست کاندیداها:} معرفی ۳۰۰ نفر با رعایت تنوع جنسیتی و تخصصی.
    \item[$\square$] \textbf{اعلام داوطلبانه دارایی‌ها:} شفافیت مالی رهبران و کاندیداهای اصلی.
    \item[$\square$] \textbf{تعهد به منشور:} امضای علنی منشور اصول بنیادین غیرقابل‌نقض.
\end{itemize}

\section{نقدهای احتمالی و پاسخ‌های طرح مهستان}

\begin{center}
\captionof{table}{تحلیل انتقادی طرح}
\begin{tabular}{R{5cm} L{9cm}}
\hline
\textbf{نقد} & \textbf{پاسخ و راهکار} \\
\hline
%\endfirsthead
%\hline
%\textbf{نقد} & \textbf{پاسخ و راهکار} \\
%\hline
%\endhead
%\hline
%\foot
«انتخابات فوری در شرایط بی‌ثباتی ممکن نیست.» & تجربه‌ی تونس نشان داد که با آمادگی جریان‌ها، انتخابات در کمتر از ۱ سال ممکن است. طرح مهستان بر آمادگی «پیش از سقوط» تأکید دارد. \\
«سیستم به جریان‌ها قدرت می‌دهد نه افراد.» & سیستم \lr{MMP} (رأی دوم فهرستی) قدرت جریان‌ها را با رأی مستقیم حوزه‌ای (رأی اول) متعادل می‌کند. \\
«دو سال برای گذار بسیار کوتاه است.» & هدف، نه پایان تمام اصلاحات، بلکه استقرار نهادهای دموکراتیک دائمی است که پس از آن مسئولیت را بر عهده می‌گیرند. \\
«خطر نفوذ عوامل نظام سابق وجود دارد.» & مکانیزم‌های \lr{Vetting} و نظارت بین‌المللی دقیقاً برای مدیریت این ریسک طراحی شده‌اند. \\
\hline
\end{tabular}
\end{center}

\begin{modernbox}{سخن پایانی با منتقدان}
هیچ طرحی بدون نقص نیست. طرح مهستان یک «پیشنهاد برای گفتگو» است. ما از هر نقد سازنده‌ای که به جای تخریب، به دنبال تقویت چارچوب‌های نهادی برای آینده‌ی ایران باشد، استقبال می‌کنیم.
\end{modernbox}
